% Part: sets-functions-relations
% Chapter: sets
% Section: russells-paradox

\documentclass[../../../include/open-logic-section]{subfiles}


\begin{document}

\olfileid{sfr}{set}{rus}
\olsection{De paradox van Russell}

Op grond van extensionaliteit mogen we de notatie
$\Setabs{x}{\phi(x)}$ gebruiken voor \emph{de} verzameling van alle
$x$ waarvoor~$\phi(x)$ geldt. Uit extensionaliteit volgt echter
\emph{werkelijk} alleen het volgende. \emph{Als} er een verzameling
bestaat met precies de objecten die aan~$\phi$ voldoen als elementen,
\emph{dan} is er maar één zo'n verzameling. Anders gezegd: als~$\phi$
is vastgelegd, is de verzameling
$\Setabs{x}{\phi(x)}$ uniek, \emph{als zij bestaat}.

Maar dit voorbehoud is belangrijk!{} Niet iedere eigenschap leent
zich voor \emph{comprehensie}. Sommige eigenschappen definiëren dus
\emph{geen} verzameling. Als alle eigenschappen dat wel deden, zouden
we regelrecht in tegenspraken belanden. Het bekendste voorbeeld
hiervan is de paradox van Russell.

Verzamelingen kunnen !!{element}s van andere verzamelingen zijn---de
machtsverzameling van een verzameling~$A$ bestaat bijvoorbeeld zelf
uit verzamelingen. Daarom is het zinvol na te gaan of een verzameling
!!a{element} van een andere verzameling is. Kan een verzameling
element van zichzelf zijn? Het begrip verzameling lijkt dat niet uit
te sluiten. Als \emph{alle} verzamelingen bijvoorbeeld samen een
collectie objecten vormen, zou men kunnen denken dat ze in één
verzameling kunnen worden samengebracht---de verzameling van alle
verzamelingen. Omdat die zelf een verzameling is, zou zij
!!a{element} van de verzameling van alle verzamelingen zijn.

De paradox van Russell ontstaat wanneer we de eigenschap beschouwen
om zichzelf niet als !!a{element} te hebben, oftewel \emph{geen
element van zichzelf te zijn}. Stel dat er een verzameling bestaat
van alle verzamelingen die zichzelf niet als !!a{element} hebben.
Bestaat
\[
R = \Setabs{x}{x \notin x}
\]
dan? We kunnen bewijzen dat dit niet zo is.

\begin{thm}[De paradox van Russell]\ollabel{thm:russells-paradox}
	Er bestaat geen verzameling $R = \Setabs{x}{x \notin x}$.
\end{thm}

\begin{proof}
Als $R = \Setabs{x}{x \notin x}$ bestaat, dan geldt
$R \in R$ dan en slechts dan als $R \notin R$. Dat is een tegenspraak.
\end{proof}

\begin{tagblock}{novice}
\begin{explain}
Laten we dit bewijs langzamer doorlopen. Als $R$ bestaat, is het
zinvol te vragen of $R \in R$ geldt. Neem eerst aan dat inderdaad $R
\in R$. De verzameling~$R$ is gedefinieerd als de verzameling van alle
verzamelingen die geen !!{element}s van zichzelf zijn. Als $R \in R$,
dan is $R$ element van zichzelf en mist die verzameling dus de eigenschap
waarmee $R$ is gedefinieerd. Maar alleen verzamelingen met die eigenschap
behoren tot~$R$. Daarom kan $R$ geen !!{element} van~$R$ zijn, dus
$R \notin R$. De verzameling~$R$ kan niet zowel wel als niet
!!a{element} van~$R$ zijn. We hebben dus een tegenspraak.

Omdat de aanname $R \in R$ tot een tegenspraak leidt, volgt $R \notin
R$. Maar ook dit leidt tot een tegenspraak!{} Want als $R \notin R$,
dan heeft $R$ zelf juist wel de eigenschap waarmee $R$ is
gedefinieerd. Net als alle andere verzamelingen die geen element van
zichzelf zijn, zou $R$ dus !!a{element} van $R$ zijn. Opnieuw kan $R$
niet zowel geen als wel !!a{element} van zichzelf zijn.
\end{explain}
\end{tagblock}

\begin{digress}
Hoe kunnen we een verzamelingenleer opzetten zonder in de paradox van
Russell terecht te komen, dus zonder de \emph{inconsistente} bewering
te doen dat $R = \Setabs{x}{x \notin x}$ bestaat? Daarvoor moeten we
axioma's formuleren die heel precies aangeven wanneer verzamelingen
bestaan (en wanneer niet).
	
De verzamelingenleer die in dit hoofdstuk wordt geschetst doet dat
niet. Ze is \emph{werkelijk na\"ief}. Ze stelt alleen dat
verzamelingen aan extensionaliteit voldoen en dat we uit gegeven
verzamelingen hun vereniging, doorsnede enzovoort kunnen vormen. Het
is mogelijk de verzamelingenleer nauwkeuriger en strenger te ontwikkelen.
\oliflabeldef{cumul:::part}{Die strengere uitwerking komt aan bod in
Deel \olref[cumul][][]{part}. Voorlopig gaan we na\"ief verder en
proberen we tegenspraken zorgvuldig te vermijden.}{}
\end{digress}
\end{document}
